\documentclass[letterpaper]{article}
\usepackage[submission]{aaai2027}
\usepackage{times}
\usepackage{helvet}
\usepackage{courier}
\usepackage{url}
\usepackage{graphicx}
\usepackage{booktabs}
\usepackage{multirow}
\usepackage{amsmath}
\usepackage{microtype}

\title{Evidence-Grounded and Repairable Research Ideation:\\
A Cross-Task Workflow with Blind Evaluation and Execution Feedback}
\author{Anonymous AAAI-27 Submission}

\begin{document}
\maketitle

\begin{abstract}
Large language models can produce plausible research proposals, but their ideas are often weakly grounded, underspecified, and difficult to evaluate or repair. We introduce a structured workflow that binds research ideation to paper evidence and baseline weaknesses, requires implementation-ready experimental plans, and closes the loop through anonymous review, mechanism-consistent repair, reference-claim verification, and execution feedback. We evaluate the workflow on three heterogeneous computer-vision research tasks under controlled generation budgets, comparing direct prompting, ResearchArena, an unrepaired focused generator, and the full system. Our protocol combines multi-model blind review with human expert judgments, component ablations, claim-verification accuracy, and cost analysis. We further present a real-data anomaly-detection case study in which execution feedback reveals a cross-category calibration failure and triggers a targeted repair. This draft deliberately treats all numerical claims as pending until the frozen AAAI-27 experiment manifest has been executed.
\end{abstract}

\section{Introduction}
Research ideation requires more than fluent proposal writing: a useful idea should identify a concrete weakness, introduce a mechanism matched to that weakness, and specify a falsifiable experiment. Existing LLM-based research agents can generate broad proposals, yet it remains difficult to determine whether additional detail reflects genuine research quality or merely verbosity. We study whether evidence constraints, structured outputs, and explicit repair checks improve research ideas under controlled conditions.

Our contributions are threefold. First, we develop an evidence-grounded, structured, and repairable ideation workflow. Second, we introduce an evaluation loop that combines anonymous pairwise review, mechanism-consistency analysis, and claim-evidence verification. Third, we connect ideation to execution feedback through a bounded real-data case study. We do not claim universal idea-generation or task-algorithm state of the art.

\section{Related Work}
\textbf{LLMs for scientific discovery and research ideation.} TODO: add verified citations for ResearchArena, AI Scientist-style systems, and controlled studies of LLM-generated research ideas.

\textbf{LLM-as-a-judge.} TODO: cover evaluator bias, position bias, verbosity bias, and human--LLM agreement.

\textbf{Evidence grounding and automated claim verification.} TODO: distinguish retrieval quality from entailment/claim support accuracy.

\section{Method}
Given a task specification, the workflow retrieves and structures paper evidence, identifies baseline weaknesses, generates focused ideas and experimental plans, obtains anonymous critiques, repairs mechanism-specific deficiencies, and verifies reference claims. A final plan can be converted into executable artifacts, whose failures are fed back into diagnosis and repair.

\subsection{Structured Evidence-Grounded Ideation}
Each idea specifies a baseline weakness, minimal new module, algorithmic objective, required scripts and data, expected tables and figures, success thresholds, and negative controls. Evidence identifiers remain traceable through generation and verification.

\subsection{Consistency-Aware Targeted Repair}
The repair stage operates on reviewer rationales and failed rubric dimensions. It must preserve idea identity and check that newly introduced losses, modules, and experiments match the proposed mechanism and task domain.

\subsection{Blind Evaluation and Claim Verification}
Ideas are formatted anonymously and evaluated in randomized A/B order. Claim verification checks paper identity and semantic support; its accuracy is evaluated against a manually annotated claim--paper subset rather than inferred from workflow pass rates.

\subsection{Execution Feedback}
Selected plans are converted into data, script, metric, and failure-check artifacts. Execution feedback may trigger calibration or implementation repair. This stage is evaluated as a case study and is not presented as an algorithmic benchmark contribution.

\section{Experimental Setup}
We use three heterogeneous tasks: physical property prediction, indoor single-image-to-3D scene generation, and industrial anomaly detection with an agent workflow. Human motion generation and 3D reconstruction remain held out. The indoor-3D evidence bank is seeded and is disclosed separately.

We compare direct prompting, a controlled ResearchArena-style baseline, focused generation without repair, a compute-matched generic self-refinement baseline, and the full workflow using the same base model, evidence access, and number of ideas. The two refinement methods use the same number of calls; all methods report actual token use and latency. Five generation replicates are planned. Evaluation includes human and multi-model blind review, position-swapped subsets, component ablations, bootstrap confidence intervals, and cost/latency accounting. Exact settings are frozen in the experiment manifest.

\section{Results}
\begin{table}[t]
\centering
\small
\caption{Main controlled comparison. All entries are pending frozen-protocol experiments.}
\label{tab:main}
\begin{tabular}{lccc}
\toprule
Method & Human win rate & Overall & Impl. readiness \\
\midrule
Direct prompting & -- & -- & -- \\
ResearchArena & -- & -- & -- \\
Focused, no repair & -- & -- & -- \\
Focused, generic refine & -- & -- & -- \\
Full workflow & -- & -- & -- \\
\bottomrule
\end{tabular}
\end{table}

\subsection{Main Comparison and Ablations}
TODO: populate only from derived tables produced from the frozen manifest. Report confidence intervals, paired effect sizes, and corrected comparisons.

\subsection{Human--LLM Judge Agreement}
TODO: report correlation, agreement, position bias, model-family bias, and parse failure rates.

\subsection{Failure, Repair, and Execution Feedback}
The physical-property case preserves a failed first repair in which mechanisms were incorrectly transferred across ideas; a second consistency-aware repair separates the mechanisms. The anomaly-detection case illustrates execution-level diagnosis, but final reporting must use a calibration/test split and include both false-positive rate and recall.

\section{Limitations and Ethical Considerations}
Idea quality is partly subjective and domain dependent. LLM judges can share biases with generators, paper retrieval may omit relevant work, and evidence verification does not establish scientific truth. Human evaluation must follow the applicable institutional requirements, disclose evaluator recruitment and compensation, and avoid collecting unnecessary personal data. Generated ideas may overlap prior work; researchers remain responsible for novelty checks, citations, experiments, and all submitted claims. Generative AI assistance used in manuscript preparation will be documented according to AAAI policy.

\section{Conclusion}
We present a framework for evidence-grounded and repairable research ideation and a protocol for evaluating it beyond fluency. The final conclusion will be calibrated to the controlled experiments and human evaluation.

\bibliographystyle{aaai2027}
\bibliography{references}
\end{document}
